\section{Using the AppImage}
\label{sec:install_appimage}

To summarize, an AppImage is a form of binary distribution in one complete package, which does not require installation of any libraries or alteration of the operating system. 
Please refer to \url{https://appimage.org/} to get additional information. \\

To obtain the COMPASS AppImage, download the latest version from \url{https://github.com/hpuhr/COMPASS/releases}. \\

Before starting, it has to be made executable (once after download) using the following command:
\begin{lstlisting}
chmod +x COMPASS-x86_64_release.AppImage
\end{lstlisting}

This is it. The application can then be run using the following command:
\begin{lstlisting}
./COMPASS-x86_64_release.AppImage
\end{lstlisting}

The following points should be considered:

\begin{itemize}  
\item The AppImage should run under any Linux distribution of similar date to Debian 10 (released July 2019) or later, but no guarantees can be made.
\item To the author's knowledge, running in virtualized operating systems is possible in some solutions, but requires additional setup (GPU acceleration).
\item Using remote desktop via VNC is slightly tricky to set up, but VirtualGL + TurboVNC works well. For a detailed setup guide, see \nameref{sec:appendix_remote_gpu} in the Appendix.
\item Geographic View rendering is performed according to the local graphics card and driver, and might be limited by their capabilities.
\end{itemize}

\section{Using the Windows Package}
\label{sec:install_windows}

The Windows package is a zip file which contains the application and all required libraries, including the Geographic View. It does not require installation. \\

To use the Windows package:

\begin{itemize}
\item Download the latest \texttt{COMPASS\_<version>\_win64.zip} from \url{https://github.com/hpuhr/COMPASS/releases}
\item Unpack the zip file into a folder of choice, it contains the folder \texttt{COMPASS\_<version>\_win64}
\item Start \texttt{compass\_client.exe} in this folder
\end{itemize}
\ \\

The following points should be considered:

\begin{itemize}
\item The Windows package was tested on Windows 10.
\item Geographic View rendering is performed according to the local graphics card and driver, and might be limited by their capabilities.
\end{itemize}
